\chapter*{Sammendrag}

Norske transportbedrifter mangler systematisk innsikt i hvordan ressursene faktisk blir brukt. Overtid håndteres reaktivt, dødtid mellom oppdrag er usynlig i dagens verktøy, og lastfordeling mellom sjåfører er avhengig av trafikkleders hukommelse. Denne masteroppgaven undersøker hvordan en algoritmestøttet ressursplanleggingsplattform kan bedre ressursutnyttelsen i norske transportbedrifter samtidig som den forblir ansvarlig overfor trafikklederen som bruker den.

Oppgaven følger paradigmet Design Science Research, anvendt gjennom Peffers' seks DSRM-aktiviteter og seks navngitte iterasjoner. Det empiriske grunnlaget er semistrukturerte intervjuer med syv trafikkledere og transportledere i norske bedrifter, analysert gjennom Braun og Clarkes seks faser for tematisk analyse. Det utviklede artefaktet, RP, er en plattform bygget på Next.js, tRPC og PostgreSQL med tre optimeringsmotorer (en grådig heuristikk, OR-Tools CP-SAT og Timefold) og et Gantt-basert dra-og-slipp-tidsskjema med menneske-i-løkken. Evalueringsrammen sammenligner motorene mot tre syntetiske datasett og kompletteres av en kravsporbarhetsmatrise.

Intervjuene avdekket et gjennomgående synlighetsgap i ressursutnyttelse, og en operatør-mot-eier-asymmetri som speiler Bainbridges \emph{Ironies of Automation}: behovet for automatisering uttrykkes av eiere og av Admmit, ikke av trafikklederne som må kjøre systemet daglig. Diskusjonen organiserer hovedfunnene under tre låste forankringskonsepter, brukt verbatim på engelsk: \textbf{Efficiency} (synlighetsgapet, asymmetrien, sammenligning av løsere), \textbf{Control} (firelagstolkningen av menneske-i-løkken anvendt på override-autoritet) og \textbf{Adaptability} (konfigurerbare myke-skranke-vekter på tvers av bedrifter). Tolv hierarkiske begrensninger (L1 til L12) avgrenser hva oppgaven kan hevde, der de viktigste er fraværet av produksjonsutrulling (L6) og brukertesting med trafikkledere (L8). Konklusjonen er at algoritmestøttet planlegging under interessentasymmetri i norsk transport er gjennomførbar når operatøren beholder autoritet over hver enkelt tildeling, systemet forklarer seg i operatørens vokabular, og det som skiller én bedrift fra en annen lever som konfigurasjon, ikke som kode.
